\section{Mechanical model and dynamics}

\subsection*{Q1: Lumped-element equation of motion}

In vacuum and without damping, the suspended mass behaves as a single
spring-mass with input force $\Fe$ and output displacement $\xm$:
\begin{equation}
  m\,\ddot{\xm} + k\,\xm \;=\; \Fe.
  \label{eq:eom-undamped}
\end{equation}
Once air damping and clamping losses are included (Q4) the equation extends to
\begin{equation}
  m\,\ddot{\xm} + b\,\dot{\xm} + k\,\xm \;=\; \Fe,
  \label{eq:eom-damped}
\end{equation}
which is the classical mass-spring-damper system, equivalent to a series RLC
circuit (Fig.~\ref{fig:rlc}).

\begin{figure}[H]
\centering
\begin{circuitikz}[scale=0.95,transform shape]
  \draw (0,0)
    to[V,l=$\Fe$,invert] (0,3)
    to[L,l=$L\equiv m$] (3,3)
    to[R,l=$R\equiv b$] (6,3)
    to[C,l=$C\equiv 1/k$] (9,3)
    to[short] (9,0)
    to[short] (0,0);
\end{circuitikz}
\caption{Series-RLC equivalent circuit of the lumped mass-spring-damper
system. Effort $\leftrightarrow$ voltage, flow $\leftrightarrow$ current.}
\label{fig:rlc}
\end{figure}

\subsection*{Q2: Fundamental resonance frequency}

\paragraph{Stiffness.}
Each of the four beams is clamped at its anchor and connected to the rigid
mass at the other end. By the symmetry of the upper/lower beam pair, the
mass-side end translates without rotation, a \emph{fixed-guided} boundary
condition. Each beam therefore contributes
\begin{equation}
  k_{b} \;=\; \frac{12\,E_{Si}\,I}{L_{b}^{3}},
  \qquad
  I \;=\; \frac{t\,W_{b}^{3}}{12},
\end{equation}
where the in-plane width $W_{b}$ is the dimension along the bending
direction. With four beams in parallel,
\begin{equation}
  k \;=\; 4 k_{b} \;=\; \frac{4\,E_{Si}\,t\,W_{b}^{3}}{L_{b}^{3}}
        \;=\; \SI{40.30}{\newton\per\meter}.
\end{equation}

\paragraph{Effective mass.}
\textcolor{blue}{Neglecting the beam mass, the moving body is the suspended
silicon drawn in Fig.~\ref{fig:schematic}: the modulation bar, the central
column, the comb-spine plate that carries the fingers (the lower ``mass''
block) and the moving comb fingers. The geometry has \emph{two}
$L\times W_{m}$ plates, the modulation bar \emph{and} the comb spine, and
$20$ moving fingers, since $N=40$ counts the comb \emph{gaps}, not the
fingers:}
\begin{equation}
  m \;=\; \rho_{Si}\,t\bigl(\,\textcolor{blue}{2}\,L\,W_{m} + L_{m}\,W_{m}
        + \textcolor{blue}{20}\,L_{c}\,W_{c}\bigr)
        \;=\; \textcolor{blue}{\SI{3.54e-12}{\kilo\gram}}.
\end{equation}
\textcolor{blue}{The first submission counted only one plate and $N=40$
fingers, giving \SI{2.33e-12}{\kilo\gram}. Adding the missing comb spine and
using the correct finger count raises the mass by $1.5\times$ and brings this
analytical estimate close to the finite-element value of Part~2.}

\paragraph{Resonance.}
\begin{equation}
  \omega_{0}=\sqrt{\frac{k}{m}}=\textcolor{blue}{\SI{3.37e6}{\radian\per\second}},
  \qquad
  f_{0}=\frac{\omega_{0}}{2\pi}\approx \textcolor{blue}{\SI{537}{\kilo\hertz}}.
\end{equation}

\subsection*{Q3: Static deflection due to gravity}

With gravity acting along the direction of motion, static balance gives
\begin{equation}
  x_{g}=\frac{m\,g}{k}
       =\frac{\textcolor{blue}{\SI{3.54e-12}{\kilo\gram}}\cdot\SI{9.81}{\meter\per\second\squared}}
              {\SI{40.30}{\newton\per\meter}}
       \approx \textcolor{blue}{\SI{0.86}{\pico\meter}}.
\end{equation}
This is more than five orders of magnitude smaller than $g_{0}$. Gravity is
irrelevant at this scale; the comb-drive force dominates by a wide margin.
